\section{Heterogeneous compute: core governance, GPU offload, and the asynchronous tower}
\label{sec:heterogeneous}

The methods in the substrate are not uniform in cost. Most return in well under a
second on a stored panel; a few -- the transfer-entropy estimators in particular --
spend their time in an $O(n^2)$ nearest-neighbour search and, run carelessly, can
saturate a workstation. This section describes how \engine{} runs that uneven load
safely. There are three concerns and one rule that ties them together. The rule is
that \emph{the number of workers must never change a number}: every heavy loop in the
engine is an order-free reduction over independent directed pairs or grid cells, so
how we split the work governs only the use of the machine, never the result. Given
that rule, the three concerns are how many CPU cores a job may use (core governance),
whether the heaviest primitive can be moved off the CPU entirely (GPU offload), and
how a job that runs for minutes rather than milliseconds is managed without blocking
the synchronous kernel (the asynchronous tower).

\subsection{Core governance: one budget, wired everywhere}
\label{sec:core-governance}

The motivating failure is concrete and was diagnosed to a single line. Three methods
-- \method{ksg\_te}, \method{ksg\_robustness}, and \method{sri\_daily} -- each
hard-coded their worker count as \code{ncores <- detectCores() - 1L}. On the
22-logical-core workstation that is $21$ \code{mclapply} workers per job, and it
ignored the job's own \code{n\_cores} request entirely. With the tower configured to
run two jobs at once and the numeric libraries left multi-threaded, two concurrent
\method{ksg\_robustness} jobs forked $2 \times 21 = 42$ worker processes, each of
which could in turn spawn its own BLAS threads on top of the fork. That multiplicative
oversubscription pinned all $22$ cores at $100\%$ and hung the machine; the reboot
left two \method{ksg\_robustness} records stranded in the \code{running} state. The
fault was not the estimator but the absence of a single place that decides how many
cores a job may use.

The fix is exactly that single place. In \file{r/\_io.R} we added one governed
budget, \method{ce\_ncores}$(p,\ \mathit{reserve})$, and wired it into all six
\code{mclapply} call sites. Its precedence is fixed and the ceiling is clamped last:

\begin{lstlisting}[language=R,caption={The single core budget (\file{r/\_io.R}).},label=lst:ce-ncores]
ce_ncores <- function(p = NULL, reserve = 2L) {
  b <- max(1L, parallel::detectCores() - as.integer(reserve))      # 1. base
  if (!is.null(p) && !is.null(p$n_cores)) {                        # 2. job override
    req <- suppressWarnings(as.integer(p$n_cores))
    if (!is.na(req) && req >= 1L) b <- req
  }
  cap <- suppressWarnings(as.integer(Sys.getenv("CE_MAX_CORES", "")))
  if (!is.na(cap) && cap >= 1L) b <- min(b, cap)                   # 3. HARD ceiling, last
  max(1L, b)
}
\end{lstlisting}

A base of \code{detectCores() - reserve} can be raised or lowered by an explicit
\code{n\_cores} parameter, but the environment ceiling \code{CE\_MAX\_CORES} is applied
afterwards and cannot be exceeded by any caller. The function always returns at least
one. The complementary half lives in the tower process manager. The job server in
\file{server/job-server.mjs} sets, for every R subprocess it spawns,
\[
  \texttt{CE\_MAX\_CORES} \;=\; \Big\lfloor \tfrac{\texttt{cpus}-2}{\texttt{MAX\_CONCURRENT}} \Big\rfloor,
\]
which evaluates to \textbf{10} on the 22-core workstation (reserving two threads for
the operating system and the Node server). Because the ceiling is divided by the
concurrency level, $\texttt{MAX\_CONCURRENT}$ jobs cannot jointly oversubscribe the
box: two concurrent jobs now use at most $20$ cores, never $42$, and a job that asks
for $21$ under a ceiling of $10$ is forced down to $10$. The same spawn pins each
numeric library to a single thread (\code{OPENBLAS\_NUM\_THREADS},
\code{OMP\_NUM\_THREADS}, \code{MKL\_NUM\_THREADS}, \code{VECLIB\_MAXIMUM\_THREADS},
\code{NUMEXPR\_NUM\_THREADS} all set to \code{1}), removing the second factor of the
oversubscription. The resulting budget is surfaced in the tower's \route{/health}
response and startup log as \code{core\_budget}, so the governing number is observable
rather than implicit.

Pinning the BLAS threads is not only a safety measure. A multi-threaded BLAS reduces
in a non-deterministic order, so summation results can differ in their last bits from
run to run; pinning each library to one thread fixes that order and makes the numeric
result reproducible. The core-governance change therefore buys a reproducibility gain
on top of the safety one, and it does so without any risk to the published numbers,
because of the order-free structure of the loops it governs.

\begin{invariant}[Worker count does not change a result]
\label{inv:worker-count}
Every loop governed by \method{ce\_ncores} is a reduction over independent units --
directed pairs for the transfer-entropy methods, $(k,\text{lag})$ cells for the
robustness grid -- whose combination (a set of per-pair records, a mean, a ranking)
does not depend on the order in which the units are evaluated or on how they are
split across workers. Capping the worker count, or pinning BLAS to one thread,
therefore changes only CPU occupancy and floating-point reduction order, never the
estimate the eval suite gates.
\end{invariant}

\subsection{GPU offload of the transfer-entropy search}
\label{sec:gpu-offload}

Core governance makes the CPU path safe; the GPU path removes the heaviest load from
the CPU altogether. The hot primitive is the Kraskov--St{\"o}gbauer--Grassberger
estimator in its Frenzel--Pompe conditional-mutual-information form
\citep{kraskov2004,frenzel2007,schreiber2000}: for each point we find $\epsilon_i$,
the max-norm distance to its $k$-th neighbour in the joint
$(\text{dst}_{t+1},\,\text{src}_t,\,\text{dst}_t)$ space, then count neighbours with
marginal max-norm distance strictly less than $\epsilon_i$ in three sub-spaces. That
search is $O(n^2)$ and is exactly the work that hung the workstation. We ported it to
two \code{numba.cuda} kernels (CUDA $12.4$, numba $0.63.1$ \citep{lam2015numba}) and
run it on an NVIDIA RTX 3000 Ada Laptop GPU (compute capability $8.9$, $8$\,GB). One
thread handles one $(\text{search},\text{point})$ pair, so a directed pair's observed
estimate and all of its surrogates run in a single launch, amortising the launch and
host--device transfer overhead. The kernels compute in FP64.

The choice of FP64 over FP32 is deliberate and is the crux of why a faster device does
not threaten the engine's discipline. The Ada GPU's FP32 path is faster still, but the
KSG estimate is the anchor that the eval suite checks; an approximate value would move
that anchor and break the gate. We therefore pay for double precision so that the GPU
result is not close to the CPU result but \emph{identical} to it. Max-norm distances
are computed by absolute value, subtraction, and maximum, all exact IEEE operations,
so the $k$-th order statistic of an identical multiset and the strict-$<$ integer
counts are the same on host and device.

That identity is enforced by a failable gate, \file{gpu/equiv\_gate.R}, which runs the
GPU helper and the engine's own CPU estimator on a small bounded case and exits
non-zero unless they agree. The observed maximum absolute difference is at most
$7\times 10^{-16}$ -- it was $6.7\times 10^{-16}$ in the NAMH re-estimation -- and the
six-decimal emitted values are bit-for-bit identical. The gate caught a subtle and
instructive discrepancy. The engine's one-dimensional neighbour count, used for the
conditioning sub-space at $\text{lag}=1$, tests membership with
\code{findInterval(v - eps)} rather than the absolute form $|dp - v| < \epsilon$. On
the rare boundary tie where $|dp - v|$ rounds to exactly $\epsilon$ but $v - \epsilon$
rounds down, the interval form \emph{includes} a point that the absolute form excludes.
The GPU helper replicates the engine's interval form host-side for that
one-dimensional case, so the published estimate (the eval anchor) is preserved exactly
rather than merely approximately. The lesson is that when a published estimator is
ported to a faster device, the gate must be the estimator's own output, including its
floating-point boundary quirks, not the textbook formula.

\begin{principle}[Gate on the estimator's output, not its textbook form]
\label{prin:gate-output}
The only quantity the eval suite gates from the transfer-entropy methods is the
observed TE (the \method{ksg\_te} band and the \method{ksg\_robustness} rankings and
Spearman correlations). Because the GPU reproduces that quantity bit-for-bit, GPU
execution can be the default without ever moving a published number. The significance
$p$-value is a different matter: it comes from IAAFT surrogates
\citep{schreiber1996}, and the engine seeds neither transfer-entropy method, so the
$p$-value is a non-reproducible draw on the CPU path and on the GPU path alike. We
never claim it reproducible, and the gate never tests it.
\end{principle}

Dispatch and fallback are handled by \method{ce\_ksg\_gpu\_pairs} in
\file{r/\_ksg\_core.R}, shared by \method{ksg\_te} and \method{ksg\_robustness}. The
default is \code{AUTO}: setting \code{CE\_GPU=0} forces the CPU path, otherwise the
dispatcher probes for a usable CUDA device, and if one is present it writes the
returns matrix as a raw column-major binary alongside a JSON spec
($n$, $k$, lag, surrogate count, the pair list), invokes \file{gpu/ksg\_gpu.py}, and
parses the per-pair results, stamping a \code{compute\_path} field. Robustness is the
point of the design: any failure at all -- \code{CE\_GPU=0}, no CUDA device, a missing
or broken helper, a malformed or short result -- returns \code{NULL}, and the caller
then runs its governed CPU \code{mclapply} path. A box with no GPU, or one with a
broken toolkit, degrades silently and safely to the path that core governance already
made safe.

\begin{figure}[t]
\centering
\begin{tikzpicture}[font=\footnotesize,
  bx/.style={draw=cegrey, thick, rounded corners=2pt, fill=celight, align=center,
    inner sep=4pt, text width=3.2cm, minimum height=8mm},
  gpu/.style={bx, fill=cegreen!14, draw=cegreen},
  cpu/.style={bx, fill=ceorange!16, draw=ceorange},
  dec/.style={draw=ceblue, thick, fill=ceblue!10, diamond, aspect=2, align=center,
    inner sep=1pt, font=\scriptsize, text width=1.9cm},
  ar/.style={-{Stealth[length=2mm]}, thick, draw=cegrey}]
  \node[bx] (entry) {\code{ce\_ksg\_gpu\_pairs()}};
  \node[dec, below=7mm of entry] (d1) {GPU available?};
  \node[gpu, below=7mm of d1] (inv) {invoke \file{ksg\_gpu.py}};
  \node[dec, below=7mm of inv] (d2) {result valid?};
  \node[gpu, below=7mm of d2] (ok) {\code{compute\_path}\\\code{="gpu"}};
  \node[cpu, right=30mm of d2] (cpu) {governed CPU\\\code{mclapply}\\(\code{"cpu"})};
  \draw[ar] (entry) -- (d1);
  \draw[ar] (d1) -- node[right,font=\scriptsize]{yes} (inv);
  \draw[ar] (inv) -- (d2);
  \draw[ar] (d2) -- node[right,font=\scriptsize]{yes} (ok);
  \draw[ar] (d1.east) -| node[near start,above,font=\scriptsize]{no} (cpu.north);
  \draw[ar] (d2.east) -- node[above,font=\scriptsize]{no} (cpu.west);
  \node[align=center, font=\scriptsize\itshape, text=cegrey, below=6mm of ok,
    text width=8.4cm] (note) {both paths return the same observed transfer entropy,
    bit-exact to $\max|\Delta| \le 7\times10^{-16}$};
\end{tikzpicture}
\caption{GPU dispatch with governed-CPU fallback. ``GPU available?'' tests
\code{CE\_GPU}$\neq$0, a present CUDA device, and a successful probe; any failure,
including a malformed result, falls back to the path that core governance already
made safe. The offload changes only where the observed transfer entropy is
computed, never its value.}
\label{fig:dispatch}
\end{figure}

The surrogate inner loop stays on the host. IAAFT is dominated by a per-surrogate FFT
rather than by the neighbour search, so the helper runs it batched over surrogates
with a batched FFT and feeds the resulting source series into the same two kernels.
The measured effect of the offload is twofold -- it is faster and it frees the CPU. On
this workstation, $12$ directed pairs with $B=99$ surrogates run in $82$ seconds at
roughly one core on the GPU against $289$ seconds across roughly seven cores on the
governed CPU, a \textbf{3.5$\times$} speed-up that also returns six cores to other
work. In the isolated trial on a deterministic coupled system the kernel search alone
runs up to \textbf{73.76$\times$} faster ($n=3000$), and between $32\times$ and
$74\times$ faster across $n \in \{2000, 3000, 5000\}$, all bit-exact. The
confirmation that matters is the full-panel re-run: the exact workload that hung the
machine -- the full G20 panel, $306$ directed pairs at $B=99$ -- completed on the GPU
in \textbf{35.6 minutes} at $100\%$ of \emph{one} core with a peak of $388$\,MB of
host memory, returning $106$ of $306$ pairs significant and the USA${\to}$Japan
estimate at $\mathbf{0.154234}$, identical to the CPU value. The saturation that hung
the 22-core box is gone, and the heaviest, most parallel primitive in the engine now
runs on the device best suited to it.

\subsection{The asynchronous tower}
\label{sec:async-tower}

A run that takes thirty-five minutes does not belong on the synchronous kernel, which
runs scale-to-zero on Cloud Run under a request timeout. Four methods are heavy enough
to need a different home: \method{ksg\_te}, \method{ksg\_robustness},
\method{namh\_pipeline}, and \method{soch\_robustness}, each flagged
\code{long\_running} in the registry. These run on the always-on workstation tower
managed by \file{server/job-server.mjs}, which executes the same parameterised-only R
registry as the kernel -- it discovers methods from \file{r/*.R} and refuses anything
else -- but asynchronously. A \route{POST /api/jobs/submit} validates the workspace
and the method name, returns a \code{job\_id} immediately, and queues the work; the
manager runs at most \code{MAX\_CONCURRENT} jobs at a time under the core budget of
Section~\ref{sec:core-governance}. Each job record and its result persist to a
\file{.jobs/} directory on disk and are reloaded on startup, so a job survives a
restart and is addressable by a permalink for thirty days; a job left
\code{running} by a crash is requeued rather than lost.

Progress is reported without compromising the engine's output contract. The kernel
requires that a method's standard output be a single pure JSON object, so progress
lines cannot simply be printed. The progress emitter \method{ce\_progress} in
\file{r/\_io.R} is therefore silent unless the environment variable \code{CE\_PROGRESS}
is set, and only the tower sets it:

\begin{lstlisting}[language=R,caption={Progress is opt-in, so the synchronous contract is never polluted (\file{r/\_io.R}).},label=lst:ce-progress]
ce_progress <- function(fraction = NA, stage = "", elapsed_s = NA) {
  if (!nzchar(Sys.getenv("CE_PROGRESS"))) return(invisible(NULL))   # silent under the kernel
  cat(toJSON(list(`__progress__` = TRUE, fraction = fraction,
                  stage = stage, elapsed_s = elapsed_s),
             auto_unbox = TRUE, na = "null"), "\n", sep = "")
  flush(stdout())
}
\end{lstlisting}

Under the synchronous kernel this is a no-op and the stdout remains a single JSON
object. Under the tower, which spawns each method with \code{CE\_PROGRESS=1}, the
method emits newline-terminated progress objects tagged \code{\_\_progress\_\_}; the
job manager parses those lines, updates the job record, and pushes them to any
subscriber over Server-Sent Events at \route{/api/jobs/:id/stream}, while the final
unterminated JSON object is read as the result. The same heavy method thus serves two
callers -- a synchronous one that sees only the answer and an asynchronous one that
sees the answer and the progress -- from one code path, with the environment, not the
source, deciding which behaviour is in force.
